\chapter{Installation}
\label{sec:installation}

\section{Requirements}

SUN-DIC has been tested on Windows, macOS, and Linux, with Linux serving as the primary development platform. The examples in this manual assume \textbf{Python 3.11}. Before proceeding, verify your Python version by running the following command from a terminal or command prompt:

\begin{lstlisting}[language=bash]
python --version
\end{lstlisting}

\noindent
\textbf{Special Installation Notes}

\begin{itemize}
	\item On Windows, the \code{ray} package may not always support the latest Python release. For example, at the time of writing, Python~3.14 was available, while \code{ray} was only supported up to Python~3.12. If installation fails because of the \code{ray} dependency, install an earlier supported version of Python.
	
	\item On Linux, the GUI requires Qt6 system libraries (for example, \code{libxcb} and \code{xcb-util}). If the GUI fails to start, these libraries may need to be installed using your system package manager.
\end{itemize}

To avoid dependency conflicts, it is strongly recommended that SUN-DIC be installed in a dedicated virtual environment. The instructions below assume that such an environment is used. Installation can be performed either with \code{conda} in an Anaconda environment or with \code{pip} in a standard \code{venv} environment.

\section{Creating a Virtual Environment}

\subsection{Using \code{conda}}

Create and activate a dedicated virtual environment named \code{sundic}:

\begin{lstlisting}[language=bash]
conda create -n sundic python=3.11
conda activate sundic
\end{lstlisting}

\subsection{Using \code{venv}}

Create a virtual environment named \code{sundic}:

\begin{lstlisting}[language=bash]
python3.11 -m venv sundic
\end{lstlisting}

\noindent
Activate the environment:

\begin{lstlisting}[language=bash]
# macOS / Linux
source sundic/bin/activate

# Windows
sundic\Scripts\activate
\end{lstlisting}

\section{Installing SUN-DIC from PyPI}

Once the virtual environment has been created and activated, install SUN-DIC using:

\begin{lstlisting}[language=bash]
pip install SUN-DIC
\end{lstlisting}

\noindent
To install the optional Jupyter notebook dependencies (required for the example notebook \code{test\_sundic.ipynb}), use:

\begin{lstlisting}[language=bash]
pip install SUN-DIC[jupyter]
\end{lstlisting}

\section{Installing from GitHub (Advanced)}

Users who wish to access the latest release, experiment with development versions, or contribute to SUN-DIC can install the software directly from the GitHub repository.  To clone and install the latest release version:

\begin{lstlisting}[language=bash]
git clone https://github.com/gventer/SUN-DIC.git
pip install ./SUN-DIC
\end{lstlisting}

\noindent
To include optional Jupyter notebook support:

\begin{lstlisting}[language=bash]
pip install ./SUN-DIC[jupyter]
\end{lstlisting}

\noindent
To obtain the latest development version, clone the \code{dev} branch instead:

\begin{lstlisting}[language=bash]
git clone --branch dev https://github.com/gventer/SUN-DIC.git
\end{lstlisting}


\section{Copying the Example Files}

The \code{copy-examples} console script is installed alongside SUN-DIC. Run it once after installation to copy the bundled example files to the current working directory:

\begin{lstlisting}[language=bash]
copy-examples
\end{lstlisting}

\noindent
This command creates the following files and directories:

\begin{itemize}
	\item \code{test\_sundic.ipynb} -- a complete worked example demonstrating the Python API within a Jupyter notebook
	\item \code{settings.ini} --- a fully documented reference configuration file for the API example, but that can also be loaded into the GUI
	\item \code{planar\_images/} --- five sample TIFF images of a planar tensile specimen for use with both the API example and the GUI
\end{itemize}

\noindent
These files provide a convenient starting point for both API and GUI users.

Optionally, the \code{copy-examples} script can also be used to copy this user manual by making use of the \code{--manual} command line flag:

\begin{lstlisting}[language=bash]
copy-examples --manual
\end{lstlisting}

\noindent
In addition to the example problem files and directories listed above, this command will also create:

\begin{itemize}
	\item \code{docs/SUN-DIC\_Manual.pdf} -- this user manual
\end{itemize}

\section{Launching the GUI}

The GUI can be launched using the following console command:

\begin{lstlisting}[language=bash]
sundic
\end{lstlisting}

\noindent
If the installation was successful, the SUN-DIC graphical interface should now open.
